\input{estilo}
\usetikzlibrary{shapes.symbols, shadows}
% La figura 1 del paper de MillWheel, con sus rótulos en inglés: cada cómputo
% son varias instancias apiladas.
\begin{document}
\begin{figura}\begin{tikzpicture}[
    comp/.style={caja, fill=white, minimum width=24mm, minimum height=13mm, align=center,
      double copy shadow={shadow xshift=1.2mm, shadow yshift=-1.2mm, fill=white, draw=tinta}},
    arista/.style={datos, draw=tinta!70},
    rot/.style={etiq, align=center}]
  \node[cloud, draw, fill=relleno, cloud puffs=10, aspect=1.6, align=center, font=\small, inner sep=1pt] (ws) at (0,0) {Web\\Search};
  \node[comp] (wc) at (4.0,0) {Window\\Counter};
  \node[comp] (mc) at (8.4,2.0) {Model\\Calculator};
  \node[comp] (sd) at (8.4,-2.0) {Spike/Dip\\Detector};
  \pic at (12.4,-2.6) {base};
  \node[etiq, align=center, anchor=north] (an) at (12.4,-2.75) {Anomaly\\Notifications};
  \draw[arista] (ws) -- node[rot, above] {Queries} (wc);
  \draw[arista] (wc.north east) -- node[rot, above left] {Window\\Counts} (mc.west);
  \draw[arista] (wc.south east) -- node[rot, below left] {Window\\Counts} (sd.west);
  \draw[arista] (mc) -- node[rot, right] {Models} (sd);
  \draw[arista] (sd) -- node[rot, above] {Anomalies} (11.85,-2.0);
\end{tikzpicture}\end{figura}
\end{document}
